% 原书物理页=印刷页5--9；1.2 全节.
\section{最初的一些例子}\label{sec:fg-1-2}

\begin{numpara}[由一个元素生成的群与循环群]\label{np:fg-1-2-1}
有理整数全体构成的集合 $\mathbb{Z}$，在加法运算下，是一个由一个元素生成的群（由 $1$ 生成），从而特别地是 Abel 群. 事实上，\emph{任何由一个元素生成的群都是 $\mathbb{Z}$ 的商群}. 的确，若 $G=\langle x\rangle$，则映射 $m\in\mathbb{Z}\mapsto x^m\in G$ 是从 $\mathbb{Z}$ 到 $G$ 上的一个满同态. 我们温习下述经典事实：
\end{numpara}

\begin{lemma*}
$\mathbb{Z}$ 的任何非零子群都是由一个元素生成的，从而形如 $n\mathbb{Z}$，其中 $n>0$ 是唯一确定的整数.
\end{lemma*}

\begin{proof}
设 $\Gamma$ 是 $\mathbb{Z}$ 的一个非零子群，并设 $n$ 为 $\Gamma$ 中最小的严格正元素. 我们要证明 $\Gamma=n\mathbb{Z}$. 设 $m\in\Gamma$ 非零. 不妨取其负元，可设 $m$ 为正. 由带余除法，存在唯一的正整数对 $(q,r)$，其中 $r<n$，使得 $m=qn+r$. 于是 $r=m-qn\in\Gamma$ 为正且 $<n$，故由 $n$ 的定义，它必为零. 从而确实有 $m=qn\in n\mathbb{Z}$.
\end{proof}

由此推知，任何循环群都同构于某个形如 $\mathbb{Z}/n\mathbb{Z}$ 的商群，即整数模 $n$ 的\newterm{fg-congruence-class}{同余类}{congruence class}构成的群. 若 $G=\langle x\rangle$，这样的同构由前述同态 $m\in\mathbb{Z}\mapsto x^m\in G$ 给出，它通过 $\mathbb{Z}/n\mathbb{Z}$ 分解，其中 $n$ 为 $x$ 的阶. 由此我们还注意到，同构 $\mathbb{Z}/n\mathbb{Z}\xrightarrow{\ \sim\ }G$ 全体与 $G$ 的\newterm{fg-generator}{生成元}{generator}全体之间有一个双射. 这便把我们引向下面这条温习：

\begin{lemma*}
$\mathbb{Z}/n\mathbb{Z}$ 的生成元全体是子集 $(\mathbb{Z}/n\mathbb{Z})^\times$，即与 $n$ 互素的整数之模 $n$ 同余类全体. 它也是环 $\mathbb{Z}/n\mathbb{Z}$ 的单位子群（即乘法下的可逆元构成的群，见下文）.
\end{lemma*}

\begin{proof}
对 $x\in\mathbb{Z}$，记 $\bar{x}$ 为它在 $\mathbb{Z}/n\mathbb{Z}$ 中的像. 设 $x$ 与 $n$ 互素. 由 Bézout 定理，存在两个整数 $a,b$ 使得 $ax+bn=1$. 于是在 $\mathbb{Z}/n\mathbb{Z}$ 中有 $a\bar{x}=\bar{1}$，这一方面表明 $\bar{x}$ 可逆且逆元为 $\bar{a}$，另一方面表明 $\langle\bar{x}\rangle\supseteq\langle\bar{1}\rangle=\mathbb{Z}/n\mathbb{Z}$. 反之，设 $\bar{x}$ 生成 $\mathbb{Z}/n\mathbb{Z}$. 则存在整数 $a$ 使得在 $\mathbb{Z}/n\mathbb{Z}$ 中 $a\bar{x}=\bar{1}$. 换言之，$ax\equiv 1\pmod n$，故 $x$ 与 $n$ 互素，且 $\bar{x}$ 可逆，其逆元为 $\bar{a}$.
\end{proof}

由上所述，映射 $\alpha\mapsto\alpha(\bar{1})$ 诱导出一个同构 $\operatorname{Aut}(\mathbb{Z}/n\mathbb{Z})\xrightarrow{\ \sim\ }(\mathbb{Z}/n\mathbb{Z})^\times$. $(\mathbb{Z}/n\mathbb{Z})^\times$ 的阶传统上记作 $\varphi(n)$，函数 $n\mapsto\varphi(n)$ 称为\newterm{fg-euler-function}{欧拉函数}{Euler's totient function}（约定 $\varphi(1)=1$）.
现在正是温习同样经典的``中国引理''\footnote{译注：即通常所说的中国剩余定理（Chinese remainder theorem）.}的时机. 在此之前，当给定两个整数 $n,m$（其中 $n$ 整除 $m$）时，我们约定把将 $\bar{1}$ 送到 $\bar{1}$ 的唯一同态称为\emph{典范同态} $\mathbb{Z}/m\mathbb{Z}\longrightarrow\mathbb{Z}/n\mathbb{Z}$，也就是说，它把整数 $x$ 的模 $m$ 同余类送到 $x$ 的模 $n$ 同余类.

\begin{lemma*}
（中国引理）设 $m_1$ 与 $m_2$ 是两个 $>0$ 且互素的整数. 那么，环的典范同态 $\mathbb{Z}/m_1m_2\mathbb{Z}\to(\mathbb{Z}/m_1\mathbb{Z})\times(\mathbb{Z}/m_2\mathbb{Z})$ 是一个同构. 特别地，这两个环的单位群同构，且 $\varphi(m_1m_2)=\varphi(m_1)\varphi(m_2)$.
\end{lemma*}

\begin{proof}
由于源与目标具有相同的基数，只需证明单射性. 而既然 $m_1$ 与 $m_2$ 被假定互素，整数 $x$ 被 $m_1m_2$ 整除当且仅当它同时被 $m_1$ 与 $m_2$ 整除.
\end{proof}

\begin{exercise*}
借助 Bézout 定理，显式地写出上述同构的逆同构.
\end{exercise*}

\begin{exercise*}
对每个素数 $p$ 与每个整数 $k>0$，证明 $\varphi(p^k)=p^{k-1}(p-1)$. 由此导出：对一切整数 $n>0$，
\[
\varphi(n)=n\prod_{\substack{p\ \text{为素数}\\ p\mid n}}\Bigl(1-\frac{1}{p}\Bigr).
\]
\end{exercise*}

由于元素的阶整除群的阶，故对每个与 $n$ 互素的整数 $x$ 有 $x^{\varphi(n)}\equiv 1\pmod n$；特别地（Fermat 小定理）：
\[
\forall\,p\ \text{为素数},\ \forall\,x\ \text{使得}\ p\nmid x:\quad x^{p-1}\equiv 1\bmod p.
\]
设 $G$ 是一个循环群，$n$ 为它的阶. 上面的第一个引理表明，$G$ 的任何子群都是循环的（且其阶整除 $n$）.

\begin{exercise*}
反之，证明对 $n$ 的每个因子 $d$，$G$ 都有一个唯一的（必为循环的）$d$ 阶子群.
\end{exercise*}

特别地，循环群（或更一般地，有限 Abel 群）是单群当且仅当它的阶是一个素数\footnote{译注：原文此处作 ``si et seulement''，按文意应为 ``si et seulement si''，今按文意译出.}.

\begin{numpara}[单位群]\label{np:fg-1-2-2}
尽管我们已经用过这一概念，这里仍回顾一下：一个（含单位的）\newterm{fg-ring}{环}{ring} $A$ 是一个 Abel 群（其运算记作加法），其上另带一个内部运算 $A\times A\to A:\ (x,y)\mapsto xy$，它满足结合律、关于运算 $+$ 满足分配律，并且具有一个记作 $1$ 的单位元. 称 $A$ 的元素 $x$ 为 $A$ 的一个\newterm{fg-unit}{单位}{unit}，如果存在 $y\in A$ 使得 $xy=yx=1$（此时这样的 $y$ 唯一，记作 $x^{-1}$）. 非零环 $A$ 的单位全体（关于乘法运算）构成一个群，记作 $A^\times$. 若 $A_1$ 与 $A_2$ 是两个环，则（像对群那样逐坐标地计算）可在乘积 $A_1\times A_2$ 上赋予一个环结构. 我们有：$(A_1\times A_2)^\times=(A_1)^\times\times(A_2)^\times$.

我们已经遇到过环 $\mathbb{Z}/n\mathbb{Z}$ 的单位群 $(\mathbb{Z}/n\mathbb{Z})^\times$. 环 $\mathbb{Z}$ 的单位群 $\mathbb{Z}^\times$ 仅由 $\{\pm1\}$ 组成；它是一个 $2$ 阶循环群. 环 $\mathbb{Z}[i]$ 的单位群是 $4$ 阶循环群. 反之，环 $\mathbb{Z}[\sqrt{2}]$ 的单位群则是无限的.
\end{numpara}

\begin{remark*}
在数论中可以证明，$\mathbb{Q}$ 的一个有限扩张的整数环之单位群是一个\newterm{fg-finitely-generated}{有限生成的}{finitely generated} Abel 群. 这样的群（非典范地）同构于它的极大有限子群与一个群 $\mathbb{Z}^r$ 的乘积. 下面的命题表明，极大有限子群总是循环的.
\end{remark*}

回顾：一个\newterm{fg-field}{域}{field}是一个非零环 $K$，满足 $K^\times=K\setminus\{0\}$. 称它是\newterm{fg-algebraically-closed}{代数闭的}{algebraically closed}，如果每个多项式 $P(X)\in K[X]$ 都\newterm{fg-split-polynomial}{分裂}{split}，即形如 $P(X)=c\prod_i(X-\alpha_i)$，其中 $\alpha_i$ 是 $K$ 中的元素. 回顾：域 $K$ 的一个\newterm{fg-algebraic-closure}{代数闭包}{algebraic closure}是一个包含 $K$ 的代数闭域 $\bar{K}$，并且它的所有元素都在 $K$ 上是代数的（即满足一个系数取自 $K$ 的多项式方程）. 这里我们承认\emph{任何域都有代数闭包}，并回顾：这样的代数闭包只是在同构意义下才唯一.
下面的结果是备考法国中学教师资格考试（agrégation）的考生们所熟知的.

\begin{proposition*}
设 $K$ 是一个域. 乘法群 $K^\times$ 的任何有限子群 $G$ 都是循环的.
\end{proposition*}

\begin{proof}
考虑多项式 $X^{|G|}-1\in K[X]$. 它至多有 $|G|$ 个不同的根. 然而，$G$ 的每个元素都是它的一个根. 因此 $G$ 恰好是它的根全体. 换言之，$G$ 就是 $|G|$ 次单位根构成的群. 于是问题化为证明这个群是循环的，即存在一个\newterm{fg-primitive-root-of-unity}{本原单位根}{primitive root of unity}（阶恰为 $|G|$ 的根）. 首先注意 $|G|$ 与 $K$ 的\newterm{fg-characteristic}{特征}{characteristic}互素（否则上述多项式将不是可分的，我们也就不会有这么多的根）.

本原 $d$ 次单位根的存在性，是对 $|G|$ 的因子 $d$ 按``整除关系''归纳地证明的. 若 $d$ 是素数，由于它异于 $\operatorname{car}(K)$，我们知道多项式 $X^d-1$ 在 $K$ 的一个代数闭包中有 $d$ 个不同的根；又因这些根都是 $X^{|G|}-1$ 的根，故它们属于 $K$. 此时，任何异于 $1$ 的 $d$ 次根都是本原的. 现在，设 $p$ 是 $|G|$ 的一个素因子，并设 $d:=|G|/p$. 假设我们已有本原 $d$ 次单位根 $x$. 那么，$X^p-x$ 在 $K$ 的一个代数闭包中的任何根都是本原 $|G|$ 次单位根（从而必属于 $G$）.
\end{proof}

把这一结果应用于一个有限域 $K$，便知 $K^\times$ 是一个循环群. 然而，要具体写出一个生成元并不容易. 趁此机会回顾域论的下列结果.

\begin{theorem}\label{thm:fg-1-2-3}
若 $K$ 是一个有限域，则它的基数是某个素数 $p$ 的幂. 反之，对每个素数 $p$ 与每个整数 $n$，存在一个基数为 $p^n$ 的域，且它在同构意义下唯一. 我们把它记作 $\mathbb{F}_{p^n}$.
\end{theorem}

\begin{proof}
在本课程中，本可以把这一结果当作已知. 下面在承认代数闭包存在性的前提下，给出论证的一个概述.

设 $K$ 是一个有限域；典范环同态 $\mathbb{Z}\longrightarrow K$ 的像是一个有限整环，故形如 $\mathbb{F}_p:=\mathbb{Z}/p\mathbb{Z}$，其中 $p$ 为素数. 这样，$K$ 典范地是一个 $\mathbb{F}_p$-向量空间，因而必是有限维的，设维数为 $d$，故它的基数为 $p^d$. 取定 $K$ 到 $\mathbb{F}_p$ 的一个代数闭包 $\overline{\mathbb{F}}_p$ 中的一个\newterm{fg-embedding}{嵌入}{embedding}. 由前一命题，$K$ 的像恰好是多项式 $X^{p^d}-X$ 的根全体. 特别地，任何其他基数为 $p^d$ 的域都同构于 $K$.

反之，固定如陈述中的 $p$ 与 $n$. 考虑 $\overline{\mathbb{F}}_p$ 中多项式 $X^{p^n}-X$ 的根全体构成的集合 $K$. 它的基数为 $p^n$. 我们证明它是一个域. 它包含 $0$ 与 $1$. 对 $x,y\in K$，有 $(x-y)^{p^n}=x^{p^n}-y^{p^n}=x-y$，故 $K$ 是一个加法子群. 对 $x,y\in K^\times$，有 $(xy^{-1})^{p^n}=xy^{-1}$，故 $K$ 的确是一个子域.
\end{proof}

\begin{example*}
$\mathbb{F}_4=\mathbb{F}_2[X]/(X^2+X+1)$.
\end{example*}

设 $K$ 是一个交换域，$n$ 是一个 $>0$ 的整数，$M_n(K)$ 是系数取自 $K$ 的 $n$ 阶方阵环. 它的单位群 $(\mathrm{M}_n(K))^\times:=\mathrm{GL}_n(K)$ 由行列式非零的矩阵组成. 更一般地，对任意交换环 $A$，系数取自 $A$ 的 $n$ 阶方阵环 $\mathrm{M}_n(A)$ 的单位群 $\mathrm{GL}_n(A)$ 由行列式 $\det(x)\in A^\times$ 的矩阵 $x$ 组成. $\mathrm{GL}_n(A)$ 对标量矩阵子群 $A^\times I_n$（它构成其中心）的商群记作 $\mathrm{PGL}_n(A)$. 映射 $\det:\mathrm{GL}_n(A)\to A^\times$ 是一个群同态，其核记作 $\mathrm{SL}_n(A)=\{x\in\mathrm{M}_n(A),\ \det(x)=1\}$. 例如，$\mathrm{SL}_n(\mathbb{Z})$ 是 $\mathrm{GL}_n(\mathbb{Z})$ 的（正规）指数 $2$ 子群. 把 $\mathrm{SL}_n(A)$ 在 $\mathrm{PGL}_n(A)$ 中的像记作 $\mathrm{PSL}_n(A)$.

\begin{remark*}
对素数的任何幂 $q$，群 $L_n(q):=\mathrm{PSL}_n(\mathbb{F}_q)$ 都是单群（$n=2$、$q=2,3$ 的情形除外）. 我们不在这里证明这一结果.
\end{remark*}

\begin{numpara}[对称群]\label{np:fg-1-2-4}
设 $X$ 是一个集合. $X$ 到其自身的双射全体（也称为 $X$ 的\newterm{fg-permutation}{置换}{permutation}）在映射的复合运算下构成一个群：这就是 $X$ 的\newterm{fg-symmetric-group}{对称群}{symmetric group} $S_X$. 现在如我们所假设的那样，设 $X$ 的基数 $n$ 有限，或等价地 $X=\{1,2,\ldots,n\}$，此时把它简单地记作 $S_n$. 它的阶为 $n!$. 其一般元素记作
\[
\sigma:\begin{pmatrix}1 & 2 & \dots & n\\ \sigma(1) & \sigma(2) & \dots & \sigma(n)\end{pmatrix}.
\]
\end{numpara}

置换 $\sigma$ 的\newterm{fg-support}{支集}{support}是指 $X$ 中被 $\sigma$ 保持固定的元素之集的补集. 容易验证：支集不相交的两个置换相互交换.

一个\newterm{fg-cycle}{轮换}{cycle}是一个在其支集上诱导出循环置换的置换 $\sigma$. 于是可以把该支集编号为 $\{i_1,\cdots,i_k\}$，使得 $\sigma(i_1)=i_2$，$\sigma(i_2)=i_3$，…，$\sigma(i_{k-1})=i_k$，$\sigma(i_k)=i_1$；此时记 $\sigma=(i_1,\ldots,i_k)$. 注意这样的编号并不唯一；两个轮换 $(i_1,\ldots,i_k)$ 与 $(j_1,\ldots,j_k)$ 相重合，当且仅当存在 $t\in\mathbb{N}$ 使得 $\forall s,\ j_s=i_{s+t}$，其中下标 $s$ 在 $\mathbb{Z}/k\mathbb{Z}$ 中取读. 支集的基数 $k$ 称为轮换的\newterm{fg-length}{长度}{length}. 我们也称之为 $k$-轮换. $2$-轮换则通常称为\newterm{fg-transposition}{对换}{transposition}.

回顾下面这个经典结果：

\begin{proposition*}
任何元素 $\sigma\in S_n$ 都可以写成
\begin{enumerate}[label=\roman*)]
\item 若干个支集两两不交的轮换之积，且这些轮换被唯一确定；
\item 若干个对换之积；
\item 若干个形如 $(1,j)$，$j=2,\cdots,n$ 的对换之积；
\item 若干个形如 $(i,i+1)$，$i=1,\cdots,n-1$ 的对换之积.
\end{enumerate}
\end{proposition*}

\begin{proof}
i) 对 $n$ 作归纳. 设 $k_1$ 为 $1\in\{1,\cdots,n\}$ 在 $\sigma$ 的各次迭代作用下的轨道的基数，并设 $c_1:=(1,\sigma(1),\cdots,\sigma^{k_1-1}(1))$，这是一个长度为 $k_1$ 的轮换. 那么 $\sigma':=c_1^{-1}\sigma$ 的支集与 $c_1$ 的支集不相交. 于是可以应用归纳假设.

ii) 验证轮换 $(i_1,\cdots,i_k)$ 等于乘积 $(i_1,i_2)(i_2,i_3)\cdots(i_{k-1},i_k)$，然后应用 i). [注意次序：我们作的是映射的复合，因此从 $(i_{k-1},i_k)$ 开始作.]

iii) 事实上，$(i,j)=(1,i)(1,j)(1,i)$.

iv) 事实上，设 $i<j$. 我们验证
\[
(j-1,j)(j-2,j-1)\cdots(i+1,i+2)(i,i+1)(i+1,i+2)\cdots(j-2,j-1)(j-1,j)=(i,j).
\]
\end{proof}

\begin{exercise*}
证明 $S_n$ 由 $(1,2)$ 与 $(1,2,\cdots,n)$ 生成.
\end{exercise*}

称 $S_n$ 中元素 $\sigma$ 的\newterm{fg-signature}{符号}{signature} $\operatorname{sgn}(\sigma)$ 为数 $\prod_{1\le i<j\le n}\frac{\sigma(i)-\sigma(j)}{i-j}=\pm 1$. 若 $\sigma$ 的符号等于 $1$（相应地，$-1$），则称 $\sigma$ 是\newterm{fg-even-permutation}{偶置换}{even permutation}（相应地，\newterm{fg-odd-permutation}{奇置换}{odd permutation}）. 例如，一个对换是奇置换；长度为 $k$ 的轮换的符号为 $(-1)^{k+1}$. 映射 $\operatorname{sgn}:S_n\to\{\pm1\}$ 是一个群同态. 它的核称为\newterm{fg-alternating-group}{交错群}{alternating group} $A_n$，其元素因而恰是偶置换. 当 $n\ge 2$ 时，它是 $S_n$ 的（正规）指数 $2$ 子群.

\begin{remark*}
只要 $n\ge 5$，群 $A_n$ 就是单群. 这是一个经典结果，将在习题课中证明.
\end{remark*}
